\input{../tex-inputs/latex-leseansicht-vorspann}

\section[Arthur Schnitzler an Gustav Schwarzkopf, 26. 4. 1928]{L04465 Arthur Schnitzler an Gustav Schwarzkopf, 26. 4. 1928}
\nopagebreak\mylabel{L04465v}
\rehead{ }\normalsize\beginnumbering\briefempfaengerindex{Schwarzkopf, Gustav@\textsc{Schwarzkopf, Gustav}!zzzSchnitzler, Arthur@\emph{von Arthur Schnitzler}!1928-04-261@{26. 4. 1928}|(be}
\toendnotes[C]{\smallbreak\pagebreak[2]}
\correspDesc{Versand  durch Arthur Schnitzler am 26. 4. 1928 in Rhodos
\newline{}Übermittlung  am 30. 4. 1928 in Trieste
\newline{}Erhalt  durch Gustav Schwarzkopf im Zeitraum [30. 4. 1928 – 3. 5. 1928?] in Wien}\toendnotes[C]{\smallbreak}
\Standort{DLA, A:Schnitzler, HS.1985.1.1897.}
\physDesc{Bildpostkarte, 321 Zeichen
\newline{}Handschrift: Bleistift, lateinische Kurrent
\newline{}Versand: Stempel: »\nobreak{}\oindex{Trieste@\textbf{Trieste}|pwk}Trieste Centro, 30. IV 28, 21–22\nobreak{}«.  }\toendnotes[C]{\smallbreak}\pstart{}{\pb}Austria\oindex{Österreich@\textbf{Österreich}|pw}\pend{}\pstart{}Hr Gustav Schwarzkopf\pend{}\pstart{}Wien I\oindex{I., Innere Stadt@\textbf{I., Innere Stadt}|pw}\pend{}\pstart{}Tiefer Graben 17\oindex{Wien@\textbf{Wien}!I., Innere Stadt@\textbf{I., Innere Stadt}!Tiefer Graben 17@\textbf{Tiefer Graben 17}|pw}\pend{}{\bigskip}
\pstart
           \noindent{}{\pb}\textcolor{gray}{\textbf{9. RODI\oindex{Rhodos@\textbf{Rhodos}|pw}. – Baluardo interno di Porta Coschino\oindex{Rotes Tor@\textbf{Rotes Tor}|pw}.}}\pend
           \vspace{1em}
\pstart
           \raggedleft{}{\pb}26. 4. 928\pend
           \vspace{0.5em}
\pstart
           eben verlassen wir Rhodus\oindex{Rhodos@\textbf{Rhodos}|pw} (das es wirklich
            gibt) – eine köstliche  Insel. Die Reise verläuft ungestört und wunderschön. Am
               29.{ }Ragusa\oindex{Dubrovnik@\textbf{Dubrovnik}|pw}, am 30.{ }Triest\oindex{Triest@\textbf{Triest}|pw}, da{\geminationn} noch ein
               {\pb}paar Tage Venedig\oindex{Venedig@\textbf{Venedig}|pw}.
            Auf baldgs Wiedersehen also und viele Grüße, auch an den verehrten Bruder\pwindex{Schwarzkopf, Max 12.\,6.\,1857 Wien – 14.\,4.\,1928 ebd.@\textsc{Schwarzkopf, Max} (12.\,6.\,1857 Wien – 14.\,4.\,1928 ebd.)|pwv}.\pend
           \pstart Herzlich Ihr
                  \spacefill\mbox{A}\pend{}\selectlanguage{ngerman}\endnumbering\briefempfaengerindex{Schwarzkopf, Gustav@\textsc{Schwarzkopf, Gustav}!zzzSchnitzler, Arthur@\emph{von Arthur Schnitzler}!1928-04-261@{26. 4. 1928}|)be}\mylabel{L04465h}
\begin{anhang}
\end{anhang}\newcommand{\dateiname}{L04465}\newcommand{\titel}{Arthur Schnitzler an Gustav Schwarzkopf, 26. 4. 1928}\newcommand{\editorInnen}{Herausgegeben von Jahnke, SelmaMüller, Martin Anton}\input{../tex-inputs/latex-leseansicht-abspann}
